\subsection{LIMITS RELATIVE TO A SUBSPACE}

Let X, Y be two topological spaces, let A be a subset of X, and let \(a \in X\) be a point of the closure of A (but not necessarily in A). Let F be the trace on A of the neighbourhood filter of a in X. If f is a mapping of A into Y, then instead of saying that \(y \in Y\) is a limit of f with respect to F and writing \(y = \lim_{\mathfrak{F}} f\) , we write

\[
y = \lim _ {x \to a, x \in \mathbf {A}} f (x)
\]

and we say that \(y\) is a limit of \(f\) at \(a\) , relative to the subspace \(A\) , or that \(f(x)\) tends to \(y\) as \(x\) tends to a while remaining in \(A\) . We have then \(y \in \overline{f(A)}\) .

If \(\mathbf{A} = \complement\{a\}\) where \(a\) is not an isolated point of \(\mathbf{X}\) , then we write \(y = \lim_{x\to a,x\neq a}f(x)\) instead of \(y = \lim_{x\to a,x\in \mathbf{A}}f(x)\) .

We make analogous definitions for cluster points.

If \(f\) is the restriction to \(\mathbf{A}\) of a mapping \(g: \mathbf{X} \to \mathbf{Y}\) , we say that \(g\) has a limit (resp. cluster point) \(y\) relative to \(\mathbf{A}\) at a point \(a \in \overline{\mathbf{A}}\) , if \(y\) is a limit (resp. cluster point) of \(f\) at \(a\) , relative to \(\mathbf{A}\) .

Let B be a subset of A and let \(a \in X\) be a point of the closure of B; if y is a limit at a, relative to A, of a map \(f: A \to Y\) , then y is also a limit of f at a, relative to B; the converse is not necessarily true. But if V is a neighbourhood in X of a point \(a \in \overline{A}\) and if f has a limit y at a, relative to \(V \cap A\) , then y is still a limit of f at a, relative to A.

Let \(a\) be a non-isolated point of \(\mathbf{X}\) , so that \(a\) is in the closure of \(\complement\{a\}\) . Then a mapping \(f: \mathbf{X} \to \mathbf{Y}\) is continuous at \(a\) if and only if we have \(f(a) = \lim_{x \to a, x \neq a} f(x)\) ; this follows immediately from the definitions.

